In fact, Hacking and Prokhorov prove that \textit{every} log terminal $\QQ$-Gorenstein degeneration of $\PP^2$ is a Manetti surface.

\begin{theorem}[{\cite[Cor.\,1.2]{HackProk}}]\label{thm:manettisurfaces}
The log-terminal $\QQ$-Gorenstein degenerations of $\PP^2$ are precisely the Manetti surfaces.
\end{theorem}

Let $M$ be a partial smoothing of a Manetti surface $M_0$ over a smooth pointed curve. If $M_t$ is a general fiber of $M$, we say a divisorial sheaf $\Oc_{M_0}(D_0)$ on $M_0$ \textit{extends to $M_t$} if there is a divisorial sheaf $\Oc_{M}(D')$ on $M$ such that $\Oc_{M}(D')|_{M_0} \cong \Oc_{M_0}(D_0)$. We call $\Oc_{M}(D')|_{M_t}$ the \textit{extension} of $\Oc_{M_0}(D_0)$ to $M_t$. The next lemma shows that any divisorial sheaf on $M_0$ that is Cartier at the smoothed singularities extends to~$M_t$. A~priori, such an extension is not unique, but once we prove that any Manetti surface has class group $\ZZ$ and that the intersection numbers are preserved in partial smoothings, it follows that the extensions are unique.

\begin{lemma}\label{extendingdivisors}
Let $M_0$ be a del Pezzo surface (a normal, log terminal surface with ample anticanonical divisor), and let $M$ be a $1$-parameter $\QQ$-Gorenstein partial smoothing of $M_0$. If $\Oc(D_0)$ is a divisorial sheaf on $M_0$ that is locally free at any partially smoothed singularities then there is a base change $T'\to T$ such that $\Oc(D_0)$ extends to a divisorial sheaf on $M \times_T T'$.
\end{lemma}

\begin{proof}
We show that the obstruction to extending $\Oc_{M_0}(D_0)$ vanishes. Assume $\Oc(D_0)$ extends to a divisorial sheaf $\Oc(D_{0,n})$ on an infinitesimal deformation $M_n$ of $M_0$ over $\CC[t]/t^n$. The obstruction to extending $\Oc(D_{0,n})$ to a divisorial sheaf on $M_{n+1}$ over $\CC[t]/t^{n+1}$ is a class in
\[
\Ext^2_{\Oc_{M_n}}(\Oc_{M_n}(D_{0,n}),\Oc_{M_0}(D_0))
\]
(see \eg \cite[\href{https://stacks.math.columbia.edu/tag/08L8}{Tag 0ECH}]{stacks-project}). The local-to-global spectral sequence shows that all the obstructions are local and supported at the non-smoothed points: \[H^2(M_n,\mathcal{H}\mathrm{om}(\Oc(D_{0,n}),\Oc(D_0))) = H^2(M_0,\Oc_{M_0}) = 0\] by Kodaira vanishing, and $\mathcal{E}\mathrm{xt}^1(\Oc(D_{0,n}),\Oc(D_0))$ is supported on the non-smoothed points). So the obstruction lives in
\[
H^0(M_n,\mathcal{E}\mathrm{xt}^2(\Oc(D_{0,n}),\Oc(D_0))).
\]
This shows the obstructions are local and supported at the non-smoothed point. But locally at this point, the deformation of $M_0$ is trivial, thus there is no local obstruction and $\Oc(D_{0,n})$ extends. Lastly, it may be necessary a priori to make a base change $T'\to T$ to make the above deformation algebraic.
\end{proof}

\begin{theorem}\label{thm:classpicofmanettisurfaces} Let $M_0$ be a $\QQ$-Gorenstein degeneration of $\PP^2$.
\begin{enumerate}
\item\label{thm:classpicofmanettisurfaces1}
$\Cl(M_0)=\ZZ$ and $\Pic(M_0)=\ZZ$. We write $\Oc_{M_0}(1)$ for the positive generator of $\Cl(M_0)$.
\item\label{thm:classpicofmanettisurfaces2}
If $A$ is the direct sum of the local class groups at the singular points of $M_0$, then the sequence:
\[
0\to \Pic(M_0)\to \Cl(M_0)\to A\to 0
\]
is exact. In particular, the map from $\Pic(M_0)$ to $\Cl(M_0)$ is multiplication by $\alpha^2=|A|$.
\item\label{thm:classpicofmanettisurfaces3}
With the notation from \eqref{thm:classpicofmanettisurfaces2}, $\Oc_{M_0}(1)$ has self intersection number $1/\alpha^2$.
\item\label{thm:classpicofmanettisurfaces4}
Let $M \to T$ be a $1$-parameter $\QQ$-Gorenstein partial smoothing of $M_0$. Let $B$ be the local class group of the smoothed singularities and let $\beta=\sqrt{|B|}\in \ZZ$. A divisor $D_0\in \Cl(M_0)$ extends to a divisor $D = D_{T'}$ on $M \times_T T'$ for a base change $T' \to T$ if and only if $\Oc_{M_0}(D_0) = \Oc_{M_0}(m \beta)$ for some integer $m$. Furthermore, over a general point, the extension $D$ is in $|\Oc_{M_t}(m)|$.
\item\label{thm:classpicofmanettisurfaces5}
If $D$, $D'$ are two $\QQ$-Cartier Weil divisors on $M$ then the intersection numbers $(D|_{M_t})\cdot (D'|_{M_t})$ are independent of $t\in T$.
\end{enumerate}
\end{theorem}

\begin{proof}
We know that $M_0$ is a partial smoothing of a weighted projective space $\PP:=\Pabc$. For weighted projective spaces the map from the class group to the local class groups is surjective. Thus it follows from Lemma \ref{extendingdivisors} that the map
\[
\Cl(M_0)\to A
\]
is surjective, which shows the sequence in (2) is exact.

The divisor $\Oc(-K_{\PP})=\Oc_{\PP}(3abc)$ extends to the divisor $K_{M_0}$ on $M_0$. Let $B$ (\resp $A$) be the local class group of the smoothed (\resp unsmoothed) singularities. By Lemma~\ref{extendingdivisors}, $\Oc_\PP(|B|)$ extends to $M_0$. Set $\beta = \sqrt{|B|}$, so $\beta=\mathrm{gcd}(|B|, 3abc$), and set $\alpha=\sqrt{|A|}$. Then $\Oc_\PP(\beta)$ extends to a divisorial sheaf $\Oc_{M_0}(1)$ on $M_0$, and $\Oc_{M_0}(1)$ generates the local class group of the unsmoothed singularities. Now we wish to prove that $\Oc_{M_0}(1)$ generates $\Cl(M_0)$.

It is known that $\Pic(M_0)=\ZZ$ (\cite[Lem.\,2.1, Prop.\,6.3]{Hacking}). Let $\Oc_{M_0}(1)\in \Cl(M_0)$ be the divisor on $M_0$ from the previous paragraph (\ie $\Oc_\PP(\beta)$ extends to $\Oc_{M_0}(1)$ on~$M_0$). We will show that $\Oc_{M_0}(1)$ generates $\Cl(M_0)$. Consider the following diagram:\vspace*{-3pt}
\[
\begin{tikzcd}
0\arrow[r] & \ZZ\cdot \Oc_{M_0}(\alpha^2)\arrow[r] \arrow[d]& \ZZ\cdot \Oc_{M_0}(1)\arrow[r]\arrow[d] & \ZZ/\alpha^2\ZZ\arrow[r]\arrow[d] &0\\
0 \arrow[r] &\Pic(M_0)\arrow[r] & \Cl(M_0)\arrow[r] & A\arrow[r] &0.
\end{tikzcd}
\]
It suffices to show that the map $\ZZ\cdot \Oc(\alpha^2)$ to $\Pic(M_0)$ is surjective.

As $K_{M_t}^2$ is constant in the family $M \to T$ and the restriction of every divisor $D|_{M_t}$ is numerically a rational multiple of $K_{M_t}$, the intersection numbers of divisors are constant in the family (proving \eqref{thm:classpicofmanettisurfaces5}). On the central fiber, if $D_0 \in |\Oc_\PP(\beta)|$, then $D_0^2 = \sfrac{\beta^2}{\alpha^2\beta^2} = \sfrac{1}{\alpha^2}$. Therefore, $\Oc_{M_0}(1)^2 = \sfrac{1}{\alpha^2}$ (proving \eqref{thm:classpicofmanettisurfaces3}).

Now, let $\Lc$ be the ample generator of $\Pic(M_0)$. Then $\Lc \equiv_{\mathrm{num}} \Oc_{M_0}(\mu)$ so $\Oc_{M_0}(\mu)\cdot \Oc_{M_0}(1)$ is an integer. Let $D_0 \in |\Oc_{M_0}(1)|$. As $\alpha^2D_0$ is Cartier, $\mu \le \alpha^2$. Therefore, because $\Oc_{M_0}(\mu)\cdot \Oc_{M_0}(1) = \mu \cdot\sfrac{1}{\alpha^2}$, we must have $\mu = \alpha^2$. This implies that $\Lc = \Oc(\alpha^2 D_0) $ and $\Oc(\alpha^2 D_0)$ generates $\Pic(M_0)$. This completes the proof of \eqref{thm:classpicofmanettisurfaces1}.

To prove \eqref{thm:classpicofmanettisurfaces4}, assume that (after possibly base changing) that $\Oc_{M_0}(\mu)$ extends to an ample generator $\Oc_{M_t}(1)$ of $\Cl(M_t)$ for a general fiber $M_t$ of $M$.
As~$H^1(M_0, \Oc_{M_0}(D_0)) = 0$, the divisor $D_0$ then extends to a divisor $D$ on $M$. By \eqref{thm:classpicofmanettisurfaces5} we have $\Oc_{M_0}(\mu)^2 = \Oc_{M_t}(1)^2$. By \eqref{thm:classpicofmanettisurfaces3} the left hand side is $\mu^2/\alpha^2$ (where $\alpha^2 = |A|$, where~$A$ is the local class group of the singularities of $M_0$). And by \eqref{thm:classpicofmanettisurfaces3} the right hand side is $\beta^2/\alpha^2$ which proves that $\Oc_{M_0}(\mu)$ extends to a divisorial sheaf on $M$ if and only if $\mu = \beta$.
\end{proof}

We thank János Kollár for suggesting improvements to the following results in this section.

\begin{theorem}\label{thm:curveinsmoothlocusdP}
Let $X_T\to T$ be a flat projective family over a smooth pointed curve $0\in T$. Assume the fibers are del Pezzo surfaces (\ie the fibers have log terminal singularities and ample anticanonical divisor) and have Picard rank~$1$. Let $\Lc$ be a divisorial sheaf on $X_T$ flat over $T$ and assume that the restriction of $\Lc$ to a general fiber is a Cartier divisor.
\begin{enumerate}
\item\label{thm:curveinsmoothlocusdP1} The dimensions of the cohomology groups $h^i(X_t,\Lc|_{X_t})$ are constant in the family.
\item\label{thm:curveinsmoothlocusdP2} Assume there is a section $s\in H^0(X_0,\Lc|_{X_0})$ such that all the singularities of~$X_0$ meeting $s=0$ are smoothed in the generic fiber of the family, and $s$ lifts to a section in $H^0(X_t,\Lc|_{X_t})$ that is reduced and irreducible. If $X_t$ is a general fiber, then the linear system $|\Lc_{X_t}|$ has at most one simple base point and a general section of $H^0(X_t,\Lc|_{X_t})$ defines a smooth curve.
\item\label{thm:curveinsmoothlocusdP3} Let $C \in |\Lc_{X_t}| $ be the smooth curve from \eqref{thm:curveinsmoothlocusdP2}. If the linear system $|\Lc_{X_t}|$ has a base point, then $-K_{X_t}\cdot C = 1$.
\end{enumerate}
\end{theorem}

\begin{proof}
The first item follows from \cite[Prop.\,2.79]{KollarModuliBook} and flatness. The second item follows because the lift of $(s = 0)$ is reduced and irreducible and contained in the smooth locus of $X_t$. Call this curve $C \subset X_t$ and note that $C$ is Gorenstein. Because $-K_{X_t}$ is ample and the Picard rank is 1, we may write $C = K_X + C + A$ where $A$ is an anticanonical divisor. By the exact sequence
\[ 0 \to \Oc_{X_t} \to \Oc_{X_t}(C) \to \Oc_C(K_C+A|_C) \to 0,\]
the linear system has base points only along $C$. To prove it is has at most one base point, as $X_t$ is del Pezzo, $H^1(X_t, \Oc_{X_t}) = 0$ so the map $\Oc_{X_t}(C) \to \Oc_C(K_C+A|_C)$ is surjective on global sections, so it suffices to study the base locus of $\Oc_C(K_C+A|_C)$. If there is a base point $p \in C$, there must be an isomorphism $H^0(C,\Oc_C(K_C+A|_C)(-p)) \cong H^0(C,\Oc_C(K_C+A|_C))$, which implies that $H^1(C,\Oc_C(K_C+A|_C)(-p)) \ne 0$. Because $C$ is Gorenstein, $\Oc(K_C)$ is dualizing, so by Serre duality we have $H^1(C,\Oc_C(K_C+A|_C)(-p)) \cong H^0(C, \Oc(-A|_C)(p))^\vee$. Because~$A$ is an ample Cartier divisor on $C$, this is nonzero if and only if $p = A|_C$, \ie there is at most one simple base point. This implies that $|\Lc_{X_t}|$ has a smooth member. Finally, because $A = -K_{X_t}$, $p = A|_C $ if and only if $\deg -K_{X_t}|_C = A|_C = 1$, which occurs only if $-K_{X_t} \cdot C = 1$.
\end{proof}

\begin{corollary}\label{cor:ampleisbpfformanetti} Let $M_0$ be the central fiber of a $\QQ$-Gorenstein degeneration of $\PP^2$.
\begin{enumerate}
\item\label{cor:ampleisbpfformanetti1} If $M$ is a $1$-parameter partial smoothing of $M_0$ and $D\subset M$ is a divisor flat over $T$, then the dimension of the the cohomology $h^i(M_t,\Oc_{M_t}(D_t))$ is constant in the family.
\item\label{cor:ampleisbpfformanetti2} If $D_0\subset M_0$ is an ample Cartier divisor, then $\Oc_{M_0}(D_0)$ is globally generated.
\end{enumerate}
\end{corollary}\label{globallygenerated}

\begin{proof}

By Theorem \ref{thm:curveinsmoothlocusdP}\eqref{thm:curveinsmoothlocusdP1}, the cohomology is constant. To prove \eqref{cor:ampleisbpfformanetti2}, note that it suffices to prove that $\Oc_{M_0}(D_0)$ is globally generated for the ample generator $D_0$ of the Picard group of $M_0$. To do this, it suffices to exhibit a section on $\Pabc$ that lifts to a section of $D_0$ contained in the smooth locus of $M_0$ by Theorem \ref{thm:curveinsmoothlocusdP}\eqref{thm:curveinsmoothlocusdP3} Such a section is necessarily reduced and irreducible as it is the generator of $\Pic(M_0)$ and is contained in the smooth locus of $M_0$, and for any Cartier divisor $D_0$ on any $\QQ$-Gorenstein degeneration $M_0$ of $\PP^2$, $-K_{M_0} \cdot D_0 \ge 3$ so Theorem \ref{thm:curveinsmoothlocusdP}\eqref{thm:curveinsmoothlocusdP3} will imply the linear system $\Oc(D_0)$ is basepoint free. To exhibit the desired sections, if $M_0 = \Pabc$, the ample generator of $\Pic(M_0)$ is $\Oc(a^2b^2c^2)$ and a general member is reduced, irreducible, and avoids the singularities of $M_0$. If $M_0$ is a partial smoothing of one of the singular points (say, the one of index $a^2$) then the divisor \hbox{$(z^{ab^2}+y^{ac^2}=0)$} avoids the unsmoothed singularities on $\Pabc$ and lifts to the ample generator of $M_0$ by Theorem \ref{thm:classpicofmanettisurfaces}\eqref{thm:classpicofmanettisurfaces2}. Finally, if $M_0$ is obtained by smoothing two of the singular points (say the ones of index $a^2$ and $b^2$), then $(z^{ab}=0)$ avoids the unsmoothed singularities on $\Pabc$ and lifts to the ample generator of $M_0$ by Theorem \ref{thm:classpicofmanettisurfaces}\eqref{thm:classpicofmanettisurfaces2}.
\end{proof}

Now we have developed the machinery needed to prove Theorem \ref{thmA:nonplanarlimits}.

\begin{definition}
We say a prime number $p$ is a \textit{Markov prime} if it appears in a Markov triple.
\end{definition}

\begin{theorem}\label{smoothlimits1}
Let $(a,b,c)\ne (1,1,1)$ be a Markov triple in non-decreasing order. If $d>2$ is a multiple of $c$ then there is a smooth limit of degree $d$ plane curves that is not planar. In particular, if $p>2$ is a Markov number that is prime, there is a nonplanar degeneration of degree $p$ plane curves.
\end{theorem}

\begin{remark}
This proves Theorem~\ref{thmA:nonplanarlimits}.
\end{remark}

\begin{proof}
First, we construct a smooth Cartier divisor on $M(c)$ of the appropriate degree, and show that it extends to a smooth divisor in the degeneration of $\PP^2$ to $M(c)$. In the special case $c=2$ (so $d$ is even, and at least 4) we can degenerate a family of even degree curves to a smooth double of a degree $d/2$ curve. Projecting from a point on the curve shows the gonality of such a curve is at most $d-2$. But the gonality of a degree $d$ plane curve is $d-1$, so we see that this double cover is not planar. So we proceed by assuming $c>2$, and we similarly show the gonality is too small.

By Theorem \ref{thm:classpicofmanettisurfaces}\eqref{thm:classpicofmanettisurfaces3}, in the specialization of $\PP^2$ to $M(c)$ the line bundle $\Oc_{\PP^2}(nc)$ specializes to the line bundle $\Oc_{M(c)}(nc^2)$. By Corollary \ref{cor:ampleisbpfformanetti}\eqref{cor:ampleisbpfformanetti2}, $\Oc_{M(c)}(nc^2)$ is globally generated, so there is a smooth divisor $C\in |\Oc_{M(c)}(nc^2)|$. By Corollary \ref{cor:ampleisbpfformanetti}\eqref{cor:ampleisbpfformanetti1}, this is a specialization of a family of curves of degree $nc$.

Now we want to show that $C$ is not a plane curve. If it were planar, it would have to have degree $nc$ and gonality $nc-1$. The idea is to show that $C$ admits a low degree pencil and therefore smaller gonality. Consider the divisorial sheaf $\Oc_{M(c)}(ab)$ on $M(c)$. Specializing $M(c)$ to $M(a,b,c)$ gives a specialization of $\Oc_{M(c)}(ab)$ to the divisorial sheaf $\Oc_{M(a,b,c)}(a^2b^2)$ on $M(a,b,c)=\Pabc$. This has two natural sections $x^{b^2}$ and $y^{a^2}$, so by Corollary \ref{cor:ampleisbpfformanetti}\eqref{cor:ampleisbpfformanetti1} we see that $h^0(\Oc_{M(c)}(ab))\ge 2$. This is the desired pencil on $M(c)$. Restricting to $C$ shows that the gonality of $C$ is at most $nab$. Now we claim that $ab< c-1$ and thus the gonality of $C$ is too small to be planar. Indeed, because $c' = 3ab - c \le \max{a,b}$, we have $c \ge 2ab$. By assumption, $c > 2$ so one of $a,b > 1$. Therefore, $c > ab -1$.
\end{proof}
